\chapter{Instalación de Linux sobre almacenamiento cifrado con LUKS}
\section{Diseño de la instalación}
\pendiente{Definir distribución y esquema de particiones. Preferir un entorno de pruebas aislado si resulta adecuado. Representar las capas de almacenamiento e identificar qué particiones quedan cifradas y cuáles no.}
\section{Instalación y configuración}
\pendiente{Realizar y documentar una instalación completa del sistema sobre una partición cifrada con LUKS. Registrar decisiones, opciones del instalador y versiones. Crear un contenedor de datos en un sistema ya instalado no demuestra por sí solo este requisito.}
\section{Verificación del arranque y del cifrado}
\pendiente{Aportar evidencias del desbloqueo al arrancar, del tipo de volumen y de la relación entre volumen cifrado y sistema de ficheros montado. Identificar cobertura de datos personales, temporales e intercambio según el diseño real.}
\section{Límites, recuperación e incidencias}
\pendiente{Explicar los límites de la protección según el estado del equipo y el alcance real de las particiones. Documentar el plan de recuperación y las pruebas realizadas sin inventar resultados. Justificar cualquier partición de arranque sin cifrar.}
